\section{An integral basis at height one}\label{h3:sec:integral-basis}

The ordinary comparison now identifies the actual coefficient map and
its mod-\(p\) cup products.  We use it to construct four maps of
\(K(1)\)-local spectra.  The first two are the unit and the determinant
class already used at height two.  The other two will be an integral
degree-minus-three class \(\delta\) and its product with that same
determinant class.  The main point is to lift the ordinary orientation
through finite Witt coefficients while keeping the transition maps and
the maps from constants visible.

Throughout this section \(p\geq5\), \(T=L_{K(3)}S\), and
\(R_1=L_{K(1)}S\).  We use the local relative Morava descent result
\Cref{h3:prop:exact-descent}, including its finite coefficient-field
extension and its naturality for the ordinary tests.  The ordinary
cochain convention remains the colimit over stable compact pole
lattices from \Cref{h3:sec:framework}.

\subsection{The two logarithms and the orientation class}

Write
\[
 G_{\mathrm f}=P_1\times\GL_2(\Z_p).
\]
The ordering of these two factors is immaterial; they are the
connected and étale frame groups on the ordinary stratum.

\begin{lemma}\label{h3:des:frame-cohomology}
For every \(\ell\geq1\), there are compatible choices of generators
for which
\[
 H^*_{\hthreects}(G_{\mathrm f},\Z/p^\ell)
       =\Lambda_{\Z/p^\ell}(h,e,d),
       \qquad |h|=|e|=1,\quad |d|=3.
\]
The degree-one generators are the two normalized logarithms,
and the restriction of \(d\) to \(\hthreeSL_2(\Z_p)\) is its orientation
generator.  Coefficient extension to \(W_\ell(k)\) gives the
same exterior algebra over \(W_\ell(k)\).  These generators are
fixed by the changes of integral frames and by coefficient
Frobenius.
\end{lemma}

\begin{proof}
The group \(\hthreeSL_2(\Z_p)\) has no element of order \(p\):
a nontrivial \(p\)-th root of unity in a two-dimensional
\(\Q_p\)-representation has minimal polynomial of degree
\(p-1>2\).  It is an analytic duality group of dimension three.
Its orientation character is trivial, since conjugation on
its trace-zero Lie algebra has determinant one.
The continuous duality theorem, in the finite completed
resolution form used in \Cref{h3:sec:coheight-one}, gives perfect
pairings between degrees \(a\) and \(3-a\) with coefficients
\(\Z/p^\ell\).

Its first cohomology is zero by the elementary-matrix argument in
\Cref{h3:ord:exterior-classes}, now with coefficients \(\Z/p^\ell\).
That argument applies at every length: choose
\(b\in\Z_p^\times\) with \(b^2-1\in\Z_p^\times\);
conjugation invariance forces a homomorphism to vanish on the upper
root subgroup because multiplication by \(b^2-1\) is surjective
on its parameter group \(\Z_p\). The same holds for the lower
root subgroup using \(b^{-2}-1\), and these two subgroups generate
\(\hthreeSL_2(\Z_p)\).
Duality gives \(H^2=0\) and identifies \(H^3\) with
\(\Z/p^\ell\).  The orientation in the completed resolution
gives these generators compatibly in \(\ell\).

The determinant extension
\[
 1\longrightarrow\hthreeSL_2(\Z_p)\longrightarrow\GL_2(\Z_p)
   \longrightarrow\Z_p^\times\longrightarrow1
\]
has the section \(a\mapsto\operatorname{diag}(a,1)\).
Its action on the orientation generator is trivial, again
by determinant one on the adjoint Lie algebra.
The finite subgroup \(\mu_{p-1}\) has exact invariants,
and the remaining \(\Z_p\) has its two-term continuous
resolution.  The Hochschild--Serre sequence thus has only
quotient columns zero and one, and kernel rows zero and three.
It gives the logarithm \(e\), a lift \(d\) of the orientation,
and their nonzero product.  Their squares are zero:
\(e^2=0\) because \(2\) is invertible and \(e\) is odd,
and \(d^2=0\) also follows from the cohomological dimension.
This proves the exterior formula for \(\GL_2(\Z_p)\).
The factor \(P_1=\Z_p^\times\) contributes the second logarithm
\(h\), by its two-term resolution.  Tensoring these finite
resolution calculations gives the asserted exterior algebra.

The constructions use determinant and oriented Lie-algebra
classes, which are preserved under conjugation and transport
of a lattice.  They have coefficients in \(\Z/p^\ell\).
Consequently coefficient Frobenius fixes them.  Extension
to the finite free coefficient ring \(W_\ell(k)\) commutes
with the finite resolution and proves the last assertion.
\end{proof}

\begin{lemma}\label{h3:des:unit-directions}
Under the natural ordinary coefficient comparison, the map
from the height-one sphere sends its degree-minus-one
mod-\(p\) generator to the connected logarithm \(h\).
Under the same comparison, the mod-\(p\) reduction of the normalized
height-three logarithm detecting \(\zeta\) has image
\[
 a h+b e,\qquad a,b\in\F_p,\quad b\ne0.
\]
Thus, writing \(\zeta\) also for its image, the ordinary
mod-\(p\) coefficient algebra is
\[
 \F_p[v_1^{\pm1}]\otimes\Lambda_{\F_p}(h,\zeta,d).
\]
In this formula the degree-minus-one \(h\)-action is the
action of the actual unit \(R_1\to L_{K(1)}T\).
\end{lemma}

\begin{proof}
We first identify the unit direction, since an abstract
cohomology isomorphism would not specify it.  In the integral \(BP_*BP\) cobar construction use the
degree-zero differential \(d=\eta_R-\eta_L\), and write bars for
reduction modulo \(p\).  The invariance of \(\bar v_1\) makes
\[
 c_{v_1}=\frac{\eta_R(v_1)-\eta_L(v_1)}p
\]
integral; \(d^2=0\) and the \(p\)-torsion-free cobar terms make it a
cocycle.  For \(0\to BP_*\xrightarrow p BP_*\to BP_*/p\to0\),
the coefficient connecting map is
\[
 \partial_{\mathrm{coeff}}(\bar v_1)=[c_{v_1}]
 \in\Ext^{1,2(p-1)}_{BP_*BP}(BP_*,BP_*).
\]
Its mod-\(p\) reduction is represented by \(\bar c_{v_1}\).
On the mod-\(p\) height-one open, \(\bar v_1\) is an invariant
unit, so \(\bar c_{v_1}/\bar v_1\) is a cocycle of cohomological
degree one and internal degree zero.  Its naturality is
literal cobar naturality: for two composable formal-group
arrows, the difference of the two images of \(v_1\) is the
sum of the two corresponding differences after transport.
Changing a trivialization by one tautological arrow
therefore changes the pulled-back cocycle by a coboundary.

Pull the height-one specialized formal law to Honda form
on the connected-marking torsor of
\Cref{h3:prop:integral-frames}.  The preceding coboundary
identity identifies \(\bar c_{v_1}/\bar v_1\) with the Honda-height-one
class.  It is nonzero: on an automorphism \(a\in1+p\Z_p\)
the integral periodic frame changes \(v_1\) by
\(a^{p-1}\), and the cocycle is
\[
 \frac{a^{p-1}-1}{p}\pmod p.
\]
At \(a=1+p\) this is \(p-1\pmod p\), a unit.
It is a unit multiple of the normalized logarithm on \(P_1\).
We choose \(h\) with the corresponding normalization.
The leading coefficient of the formal-group arrow has
been retained here through the periodic frame \(U\);
discarding that frame would discard this calculation.

The height-one sphere calculation identifies its mod-\(p\)
groups with
\(\F_p[v_1^{\pm1}]\otimes\Lambda(h)\).
For the Adams self-map normalized in
\Cref{h3:prop:bounded-pole-finiteness}, the image of the Moore bottom
cell has \(BP\)-Hurewicz coefficient \(\bar v_1\).  Its Moore boundary
is a unit multiple of \(\alpha_1\).  The normalized cobar terms for
\(BP\) are even and \(p\)-torsion-free, so the comparison in
\Cref{h3:sec:descent} between the termwise cone and the cone of
totalization applies to this cofiber sequence.
For the lift \(v_1\), the coefficient connecting cocycle is
\(c_{v_1}\), while the actual desuspended boundary has leading
detector \(-[c_{v_1}]\).  The internal degree \(2(p-1)\) is
even, and strict desuspension contributes no further sign.
Choose the negative of this boundary class, map it to the Moore
quotient, and invert the same self-map.  The resulting source
generator in degree \(-1\) has detector
\(\bar c_{v_1}/\bar v_1\), the logarithm normalized above.
The odd-prime Adams-summand fiber is realized
with the unit as its first map in \cite[Section~4.2, pp.~30--32]{BB}.
Its logarithm, mod-\(p\) algebra, and unit-induced rational splitting
are given in \cite[Section~4.3, Remarks~4.16--4.17 and (4.18)]{BB}.
The maps of the preceding paragraph are the unit maps of
the formal-group cobar construction.  Hence their \(h\)
is precisely the \(R_1\)-action, rather than an independently
chosen degree-one direction.

For the second direction use the actual determinant
identity in \Cref{h3:prop:integral-frames}: the determinant
of the middle height-three frame is the product of the
norm of the connected height-one frame and the determinant
of the rank-two kernel frame.  It is an identity of
integral determinant lattices.  Applying the normalized
logarithm yields the sum of the two logarithms, with
the signs introduced by inverse-frame conventions.
In particular its coefficient on \(e\) is a unit modulo
\(p\).  This is the class \(\zeta\), by its determinant
construction.  The change of basis from \((h,e)\) to
\((h,\zeta)\) is invertible.

Finally, \Cref{h3:thm:ordinary-comparison,h3:des:frame-cohomology}
identify the actual coefficient algebra with this exterior
algebra.  The comparisons preserve products and the
two unit maps, proving all assertions.
\end{proof}

The coefficient calculation now has columns only from zero
to five.  Since the internal degrees are still multiples
of \(2(p-1)\), no mod-\(p\) differential is possible.
Cup products give the associated-graded multiplication.
For the actions used here one may use a Moore pairing
and the height-one Adams self-map; their cobar actions
are the same cup actions.  We next lift the degree-three
orientation before using these products to construct an
equivalence of spectra.

\subsection{Finite Witt length and ordinary integral coefficients}

Put
\[
 A_\ell=(W(k)/p^\ell)[[x,y]],\qquad
 D_\ell=A_\ell[x^{-1}].
\]
Write \(\phi_p(a)=a^p\) for coefficient \(p\)-Frobenius on the
characteristic-\(p\) rings \(B_0\) and \(B_{\hthreepk}\), and let
\(\sigma_k\) be the Witt automorphism of \(W(k)\) lifting
\(a\mapsto a^p\) on \(k\).  This notation is distinct from the
chosen \(q\)-power operator used in the ordinary root comparison.
The symbol \(D_\ell\) is separate from the distribution parameter
of \Cref{h3:sec:coheight-one}.
The action of \(P_3\) preserves \(A_\ell\) and satisfies
\(g(x)=xu_g+ph_g\), with \(u_g\) a unit and \(h_g\in A_\ell\).
It therefore extends to \(D_\ell\), since the second term
is nilpotent relative to the invertible first term.
An action on the uncompleted mixed-characteristic
localization \(W(k)[[x,y]][x^{-1}]\) is not needed.

To specify the cochain topologies at finite length, use
the following stable compact additive lattices, for \(d\geq0\):
\begin{equation}\label{h3:des:witt-pole-lattices}
 \mathcal L_{\ell,d}=[x]^{-d}W_\ell(A),\qquad
 \mathcal M_{\ell,d}
       =\sum_{j=0}^{\ell-1}p^j x^{-d-j}A_\ell.
\end{equation}
The first has its \(i\)-th Witt coordinate in
\(x^{-dp^i}A\).  These lattices exhaust \(W_\ell(B_0)\);
the bounds are cofinal with finite coordinate pole bounds.
Moreover,
\[
 x^{-d}A_\ell\subset\mathcal M_{\ell,d}
       \subset x^{-d-\ell+1}A_\ell,
\]
so the second family is cofinal with bounded poles in
\(D_\ell\).  Stability of the first follows from
\(g([x])=[x][\overline u_g]\).
For the second, expand \((xu_g+ph_g)^{-d-j}\)
modulo \(p^{\ell-j}\); a term with an extra \(p^a\)
has pole at most \(d+j+a\), as required by the
\((j+a)\)-th summand.  These actions are continuous
in the compact lattice topologies.
All cochains on either union mean the filtered colimit
of the continuous cochains on these stable lattices.
In particular, no coordinatewise pole-bound set that
fails to be additive is being used as a coefficient module.

\begin{lemma}\label{h3:lem:ghost-comparison}
For each \(\ell\geq1\), the last-ghost formula
\begin{equation}\label{h3:des:ghost-formula}
 \gamma_\ell(a_0,\ldots,a_{\ell-1})
   =\sum_{i=0}^{\ell-1}
       p^i\widetilde a_i^{\,p^{\ell-1-i}}\quad\hbox{in }D_\ell
\end{equation}
for arbitrary lifts \(\widetilde a_i\in D_\ell\) defines
a natural equivariant unital ring map
\[
 W_\ell(B_0)\xrightarrow{\gamma_\ell}D_\ell.
\]
Both this map and the natural coefficient map
\[
 \beta_\ell:W_\ell(B_0)\longrightarrow W_\ell(B_{\hthreepk})
\]
induce isomorphisms on \(P_3\)-cohomology with the prescribed
cochain conventions.  For \(\ell\geq2\), let
\(R_\ell:W_\ell(B_0)\to W_{\ell-1}(B_0)\) be Witt truncation and
let \(\operatorname{red}_{\ell,\ell-1}:D_\ell\to D_{\ell-1}\)
be coefficient reduction.  Then
\begin{equation}\label{h3:des:ghost-transition}
 \operatorname{red}_{\ell,\ell-1}\gamma_\ell
     =\gamma_{\ell-1}R_\ell W_\ell(\phi_p)
       \qquad(\ell\geq2).
\end{equation}
The map \(\gamma_\ell\) fixes the constant
\(\Z/p^\ell=W_\ell(\F_p)\).
\end{lemma}

\begin{proof}
The binomial theorem gives
\[
 a\equiv b\pmod p
 \quad\Longrightarrow\quad
 a^{p^j}\equiv b^{p^j}\pmod {p^{j+1}}
 \qquad(j\geq0).
\]
For the \(i\)-th summand of \eqref{h3:des:ghost-formula},
changing a lift consequently changes the result by a
multiple of \(p^\ell\).  This proves independence of all
lifts.
Lift two Witt vectors to \(W_\ell(D_\ell)\), and use their
Witt sum or product to lift the corresponding operation
over \(B_0\).  The last ghost component on
\(W_\ell(D_\ell)\) is a ring homomorphism.  Its reduction
and the lift independence prove the ring identities for
\(\gamma_\ell\).  The construction is natural for maps
of the reduction \(D_\ell\to B_0\), and hence equivariant
for the actual \(P_3\)-action on \(D_\ell\).
These are the elementary ghost identities of
\cite[Lemma~1 and Proposition~2]{Hesselholt}.

The map is continuous on each compact pole lattice.
Indeed, coefficientwise lifts and the finite polynomial
\eqref{h3:des:ghost-formula} prove continuity on \(W_\ell(A)\);
its ring property then gives
\[
 \gamma_\ell(\mathcal L_{\ell,d})
       \subset x^{-dp^{\ell-1}}A_\ell
       \subset\mathcal M_{\ell,dp^{\ell-1}}.
\]
It therefore acts on the specified cochain colimits.

Filter the source by \(V^iW_{\ell-i}(B_0)\) and the target
by \(p^iD_\ell\), for \(0\leq i\leq\ell\).
The source filtration iterates the one-coordinate
restriction--Verschiebung sequence in
\cite[proof of Lemma~9, p.~6]{Hesselholt}.
Their \(i\)-th graded pieces are additively \(B_0\);
the induced map is
\[
 a\longmapsto a^{p^{\ell-1-i}}.
\]
The quotient sequences remain exact on cochains.
For a Witt quotient the section is the coordinate
section \(a\mapsto V^i[a]\); for the corresponding
target quotient use \(a\mapsto p^i\widetilde a\) with
coefficientwise lifts.  These sections are continuous
and send a compact bounded-pole family into some
bounded-pole stage of \eqref{h3:des:witt-pole-lattices}.
They need not be additive or equivariant: their role
is to give degreewise surjectivity of continuous
cochains.  The kernels and differentials are the
equivariant ones.

By \Cref{h3:thm:ordinary-comparison}, the coefficient
\(p\)-Frobenius \(\phi_p\) is an isomorphism on
\(H^*(P_3,B_0)\).
The long exact cohomology sequences and induction on
the finite filtration now prove the assertion for
\(\gamma_\ell\).
For \(\beta_\ell\), use the two Verschiebung
filtrations.  Their graded maps are the natural
comparison \(B_0\to B_{\hthreepk}\), which is a cohomology
isomorphism by the same theorem.  Witt-coordinate
sections on the analytic side are continuous, so the
identical finite-filtration argument proves the
assertion for \(\beta_\ell\).

In characteristic \(p\), truncated Witt Frobenius is
\[
 R_\ell W_\ell(\phi_p)(a_0,\ldots,a_{\ell-1})
          =(a_0^p,\ldots,a_{\ell-2}^p);
\]
see \cite[Lemma~8]{Hesselholt}.
Substitution into \eqref{h3:des:ghost-formula} proves
\eqref{h3:des:ghost-transition}.  Since \(\gamma_\ell\)
is unital, it is the identity on \(\Z/p^\ell\).
Its restriction to \(W_\ell(k)\), under the identification
with \(W(k)/p^\ell\), is \(\sigma_k^{\ell-1}\);
we do not identify that restriction with the identity.
\end{proof}

\begin{lemma}\label{h3:des:witt-constants}
The actual ordinary constants map gives natural equivalences
\[
 \eta_\ell:
 R\Gamma_{\hthreects}(G_{\mathrm f},W_\ell(k))
    \xrightarrow{\ \simeq\ }
 R\Gamma_{\hthreects}(P_3,W_\ell(B_{\hthreepk})).
\]
They commute with truncation, coefficient Frobenius, and
the maps from \(\Z/p^\ell\) constants.  The maps
\(\beta_\ell,\gamma_\ell\) in \Cref{h3:lem:ghost-comparison}
preserve these latter constant maps as well.
\end{lemma}

\begin{proof}
On the ordinary quotient stack, the comparison of
\Cref{h3:prop:ordinary-constants} is induced by its map to
\(BG_{\mathrm f}\) and by the coefficient unit
\(k\to\mathcal O^\flat\).  Apply \(W_\ell\) to the
coefficient sheaves on this same diagram.  It gives
the displayed map before taking cohomology, and all
the asserted naturalities follow from the natural
Witt operations.  Its Verschiebung filtration has
the length-one constants comparison on each graded
piece.  The continuous coordinate sections make
the associated short exact sequences exact on the
computing cochains, including their profinite
parameters.  Finite-filtration induction from
\Cref{h3:prop:ordinary-constants} proves the equivalence.

The coefficient map \(\beta_\ell\) preserves constants
by definition.  For \(\gamma_\ell\), the composite
\[
 C^*_{\hthreects}(P_3,\Z/p^\ell)
   \longrightarrow C^*_{\hthreects}(P_3,W_\ell(B_0))
   \longrightarrow C^*_{\hthreects}(P_3,D_\ell)
\]
is exactly the ordinary coefficient inclusion, since
\(\gamma_\ell\) is the identity on \(\Z/p^\ell\).
Thus these are statements about the specified maps
from constants, not about an abstract isomorphism
of their images in cohomology.
\end{proof}

\begin{lemma}\label{h3:des:finite-length-model}
The ordinary spectrum \(L_1(T/p^\ell)\) is the totalization
of the termwise \(L_1(-/p^\ell)\) Morava descent object.
Its terms are even, and their internal degree-zero
cochain complex is the \(D_\ell\) complex above.
The corresponding spectral sequence is strongly
convergent, has columns \(0\leq s\leq5\), and has
no differentials.  All its abutment groups are finite
in each stem.  In particular there is a natural
identification
\begin{equation}\label{h3:des:minus-three-edge}
 \pi_{-3}L_1(T/p^\ell)
       \cong H^3_{\hthreects}(G_3(k),D_\ell),
\end{equation}
where the right side includes the coefficient
Galois descent just described.
\end{lemma}

\begin{proof}
The totalization assertion is
\Cref{h3:prop:exact-descent}.
Put \(M_\ell=\mathbb S_{(p)}/p^\ell\) and \(E=E_3(k)\).
The completion map identifies \(M_\ell\) with \(S/p^\ell\),
naturally in the ordinary Moore quotient transitions: its cofiber
has zero mod-\(p\) quotient, so multiplication by \(p\) is an
equivalence there and every mod-\(p^\ell\) quotient vanishes.
Each finite \(M_\ell\) has type one, since its rational homology
vanishes and its \(K(1)\)-homology is nonzero.
Let \(\psi:\Sigma^{2(p-1)}M_1\to M_1\) be the normalized Adams
\(v_1\)-map chosen in the proof of
\Cref{h3:prop:bounded-pole-finiteness}; its action on \(E/p\)
is multiplication by \(xU^{p-1}\).
The periodicity theorem
\cite[Theorem~1.5.4(i)--(ii), pp.~9--10]{RavenelNilpotence}
supplies a \(v_1\)-map \(f_\ell:\Sigma^{d_\ell}M_\ell\to M_\ell\)
and, for the ordinary reduction \(r_\ell:M_\ell\to M_1\),
positive integers \(i_\ell,j_\ell\) such that
\(d_\ell i_\ell=2j_\ell(p-1)\) and
\[
 r_\ell\circ f_\ell^{i_\ell}
   \simeq \psi^{j_\ell}\circ
          \Sigma^{2j_\ell(p-1)}r_\ell.
\]
Here the theorem is applied after a common suspension to finite
CW-complexes and then desuspended to spectra.  Multiplication of
the Adams map by the integer unit used in its normalization
preserves its \(v_1\)-map property.  Fix this iterate
\(\varphi_\ell=f_\ell^{i_\ell}\), of degree
\(e_\ell=2j_\ell(p-1)\).

The cofiber sequence \(E\xrightarrow{p^\ell}E\to E/p^\ell\)
and evenness of \(E/p^\ell\) give an isomorphism
\[
 [\Sigma^{e_\ell}E/p^\ell,E/p^\ell]_E
       \xrightarrow{\ \simeq\ }\pi_{e_\ell}(E/p^\ell)
\]
by evaluation on the unit: the possible kernel is a quotient of
\(\pi_{e_\ell+1}(E/p^\ell)=0\), and multiplication by
\(p^\ell\) vanishes on the group on the right.
Thus the \(E\)-module map induced by \(\varphi_\ell\) is determined
by a scalar \(\theta_\ell\) in this group.  Applying \(E\wedge-\)
to the chosen reduction square proves
\[
 \theta_\ell\bmod p=(xU^{p-1})^{j_\ell},\qquad
 \theta_\ell=U^{j_\ell(p-1)}c_\ell,\qquad
 c_\ell=x^{j_\ell}+p b_\ell\quad(b_\ell\in A_\ell).
\]
A lift of \(\theta_\ell\) to \(\pi_{e_\ell}E\) gives the same
\(E\)-module map by the evaluation isomorphism.  Tensoring with an
ordinary \(E\)-module \(N_Q\) therefore makes
\(\varphi_\ell\) act on \(\pi_*(N_Q/p^\ell)\) by multiplication
by this scalar.

The two localizations of \(A_\ell\) at \(c_\ell\) and at \(x\)
agree canonically.  Indeed, in \(A_\ell[x^{-1}]\) the factorization
\(c_\ell=x^{j_\ell}(1+p b_\ell x^{-j_\ell})\) makes \(c_\ell\)
invertible.  In \(A_\ell[c_\ell^{-1}]\), the factorization
\(x^{j_\ell}=c_\ell(1-p b_\ell c_\ell^{-1})\) makes \(x\)
invertible.  Both parenthesized factors have finite geometric-series
inverses because \(p^\ell=0\).
Miller's odd-prime result for \(M_1\)
\cite[Theorem~4.11 and Corollary~4.12]{Miller1981}, together with
the thick-subcategory argument in
\cite[Section~7.5, after Conjecture~7.5.5, p.~85]{RavenelNilpotence},
identifies the telescope of a \(v_1\)-map with \(L_1M_\ell\)
for every finite type-one \(M_\ell\).  Taking the chosen iterate
does not change the telescope, since it selects a cofinal subsequence.
As \(L_1\) is smashing, this is an equivalence under \(N_Q/p^\ell\)
between \(L_1(N_Q/p^\ell)\) and the telescope of its
\(\varphi_\ell\)-action.

For the actual term \(N_Q=\mathcal N_3^q\), with \(Q=P_3^q\),
\Cref{h3:des:flat-functions} gives even, \(p\)-torsion-free
coefficients.  Coefficientwise continuous lifts identify
\(\pi_0(N_Q/p^\ell)=C_{\hthreects}(Q,A_\ell)\).
Homotopy groups commute with the telescope, so the scalar calculation
and the two localizations above give
\[
 \pi_0L_1(N_Q/p^\ell)
   \cong C_{\hthreects}(Q,A_\ell)[x^{-1}]
   =\colim_d x^{-d}C_{\hthreects}(Q,A_\ell).
\]
The odd groups vanish, and the other even groups are its periodic
translates.  This is the algebraic bounded-pole localization;
the stable cofinal lattices \(\mathcal M_{\ell,d}\) specify it in
\eqref{h3:des:witt-pole-lattices}.

These coefficient identifications preserve the actual transition
maps.  Each is the unique extension of the map on coefficients
induced by the \(L_1\)-localization unit, by the universal property
of localization of a module at \(x\).  The ordinary Moore quotient
transition is induced by the cofiber square with vertical maps \(p\)
and \(\hthreeid\); it reduces coefficients and preserves the fixed
coordinate \(x\).  Naturality of the localization unit therefore
identifies its localized coefficient map with ordinary reduction
\(D_\ell\to D_{\ell-1}\).  The same argument applies to every
coface and codegeneracy of the actual descent object: coefficient
actions send \(x\) to \(xu_g+ph_g\), which is invertible in
\(D_\ell\), and precomposition on continuous functions preserves
the bounded-pole colimit.  Hence the degree-zero coefficient
complex is the stated \(D_\ell\) complex, with its actual cochain
maps, ordinary Moore transitions, and maps from constants.

The tame central units kill internal degrees
not divisible by \(2(p-1)\).
For all remaining degrees the finite \(p\)-filtration
has associated graded the mod-\(p\) coefficient
complex, trivialized by powers of \(v_1\).
By \Cref{h3:thm:ordinary-comparison,h3:des:frame-cohomology}
that complex has finite cohomology in columns
zero through five, and vanishes above five.
The long exact sequences of the finite
\(p\)-filtration give those same finiteness and
vanishing conclusions in every length.
There is no assertion that \(v_1\) itself is an
integral invariant in length \(\ell\); only the
associated-graded trivialization is used.

A differential has \(2(p-1)\mid r-1\), so the first
candidate has \(r=2p-1\geq9\).  It cannot fit the
column range \(0\leq s\leq5\).
Thus the spectral sequence collapses, and its
finite column range proves degreewise finiteness
of the abutment.  Strong convergence comes from
\Cref{h3:prop:exact-descent}.

Finally, in stem \(-3\) one must have \(t-s=-3\)
with \(0\leq s\leq5\) and \(t\in2(p-1)\Z\).
The unique solution is \((s,t)=(3,0)\).
There is consequently no additive extension or
choice of filtration quotient in
\eqref{h3:des:minus-three-edge}.  The construction is
natural in the Moore quotient maps and in
coefficient constants.  This proves its stated
naturality as well.
\end{proof}

The finite coefficient calculation has now isolated one degree-three
line at every length, with ordinary reduction as the spectral
transition.  To turn a compatible family on these lines into a map of
spectra, we spell out the completion step.

\begin{lemma}[Completion inside height one]\label{h3:lem:height-one-completion}
Let \(M\) be an \(E(1)\)-local \(S\)-module.  The natural map
\[
 M\longrightarrow \widehat M_p:=
       \operatorname*{holim}_{\ell\geq1} M/p^\ell
\]
is its \(K(1)\)-localization in \(S\)-modules.  Here
\(M/p^\ell\) is the cofiber of multiplication by \(p^\ell\),
with the ordinary Moore quotient transitions.
\end{lemma}

\begin{proof}
Write \(S/p^\infty=\operatorname*{colim}_\ell S/p^\ell\), with
the usual inclusions of Moore spectra.  Duality of the finite Moore
\(S\)-modules identifies
\(M/p^\ell\simeq F_S(\Sigma^{-1}S/p^\ell,M)\).  The inclusion
\(S/p^\ell\to S/p^{\ell+1}\) is induced by the square with vertical
maps \(\hthreeid\) and \(p\); dualizing gives the ordinary quotient
transition induced by \(p\) and \(\hthreeid\).  Thus these
identifications respect the tower.  Applying the internal function functor to
\[
 \Sigma^{-1}S/p^\infty\longrightarrow S\longrightarrow S[1/p]
\]
therefore gives a natural fiber sequence
\[
 F_S(S[1/p],M)\longrightarrow M\longrightarrow\widehat M_p.
\]
Multiplication by \(p\) is invertible on the first term.  Moreover
\[
 F_S(S[1/p],\widehat M_p)
 \simeq F_S(S[1/p]\otimes_S\Sigma^{-1}S/p^\infty,M)=0.
\]
Thus \(\widehat M_p\) is local with respect to the Moore quotient:
maps from any \(p\)-invertible \(S\)-module vanish, since such a
module is an \(S[1/p]\)-module and the displayed function object is
zero.  Each \(M/p^\ell\) is \(E(1)\)-local, and local objects are
closed under limits, so \(\widehat M_p\) is also \(E(1)\)-local.

In the \(E(1)\)-local category the Moore-acyclic objects are exactly
the rational objects: the Moore condition makes \(p\) invertible,
and conversely rationalization kills every Moore spectrum.  The
height-one fracture square identifies these with the \(K(1)\)-acyclic
\(E(1)\)-local objects.  Hence Moore localization and
\(K(1)\)-localization have the same acyclic objects and the same
local objects in this category.  The \(E(1)\)-localization unit is
itself a \(K(1)\)-equivalence, so passing first to \(E(1)\)-local
objects does not change \(K(1)\)-localization.  The displayed natural
map is therefore the asserted localization in \(S\)-modules.
\end{proof}

\begin{proposition}\label{h3:prop:integral-delta}
There is a class
\[
 \delta\in\pi_{-3}L_{K(1)}T
\]
whose reduction is the orientation generator \(d\)
in \Cref{h3:des:unit-directions}.
Moreover,
\[
 \pi_{-3}L_{K(1)}T\cong\Z_p,
\]
with \(\delta\) a generator after the chosen orientation
normalization.  The finite-length identifications used
to construct it preserve the given maps from
\(P_3\)-cochains with integral constant coefficients.
\end{proposition}

\begin{proof}
Choose the compatible classes
\(o_\ell\in H^3(G_{\mathrm f},\Z/p^\ell)\)
of \Cref{h3:des:frame-cohomology}.
Write \(\rho_\ell\) for the cohomology isomorphism
induced by \(\beta_\ell\), and set
\[
 \nu_\ell=\rho_\ell^{-1}\eta_\ell(o_\ell),\qquad
 z_\ell=\gamma_\ell(\nu_\ell)
       \in H^3(P_3,D_\ell).
\]
Here the constant coefficient extension to \(W_\ell(k)\)
is understood.  The naturalities in
\Cref{h3:des:witt-constants} give
\(R_\ell\nu_\ell=\nu_{\ell-1}\).
Coefficient Frobenius fixes \(o_\ell\), because it
fixes \(\Z/p^\ell\); injectivity of \(\rho_\ell\)
then gives
\[
 W_\ell(\phi_p)(\nu_\ell)=\nu_\ell.
\]
Using \eqref{h3:des:ghost-transition} we obtain
\[
 \operatorname{red}_{\ell,\ell-1}(z_\ell)
   =\gamma_{\ell-1}R_\ell
          W_\ell(\phi_p)(\nu_\ell)
   =z_{\ell-1}.
\]
Thus the Frobenius twist in the ghost transition
does not obstruct compatibility of these actual
integral orientation classes.

By \Cref{h3:des:frame-cohomology,h3:des:witt-constants,h3:lem:ghost-comparison}, the degree-three cohomology
over \(k\) is one free copy of \(W_\ell(k)\).
The determinant and orientation constructions are
preserved under coefficient Galois action and
transport of the integral frames.  Normal-basis
descent therefore identifies its descended group
with \(\Z/p^\ell\).  The elements \(z_\ell\) are
generators and their transitions are ordinary
reduction.  This proves the corresponding assertions
for \(\pi_{-3}L_1(T/p^\ell)\) by
\eqref{h3:des:minus-three-edge}.

Apply \Cref{h3:lem:height-one-completion} to \(L_1T\).  Exactness of
\(L_1\) identifies its Moore quotients with \(L_1(T/p^\ell)\).
Consequently
\[
 L_{K(1)}T\simeq
       \operatorname*{holim}_\ell L_1(T/p^\ell).
\]
All finite-length homotopy groups are finite by
\Cref{h3:des:finite-length-model}; their inverse systems
are Mittag--Leffler, so the Milnor
\(\lim^1\)-term vanishes in every stem.
The compatible generators \(z_\ell\) therefore
give an actual class \(\delta\), and their inverse
limit is \(\Z_p\).
The last assertion follows from the equality of
the constant cochain maps in
\Cref{h3:des:witt-constants} and from the natural
edge identification in
\Cref{h3:des:finite-length-model}.
\end{proof}

\begin{remark}\label{h3:des:constants-order-bridge}
The preceding argument does not identify the usual
full Witt tower on the uncompleted \(B_0\) with the
ordinary \(D_\ell\) tower.  For later use, the exact
statement needed for any class from integral
\(P_3\)-constants is simpler.
Its image in \(W_\ell(B_0)\) is fixed by coefficient
Frobenius, so \eqref{h3:des:ghost-transition} gives an
ordinary reduction-compatible family in \(D_\ell\).
The two cohomology isomorphisms in
\Cref{h3:lem:ghost-comparison} preserve the order of
each finite-length element.
Hence unbounded \(p\)-power order of the corresponding
analytic Witt images implies unbounded order of
these specified ordinary images.  Combined with
\Cref{h3:prop:integral-delta}, this is the bridge used
by the rational detector in \Cref{h3:sec:rational}.
\end{remark}

\subsection{The actual height-one equivalence}

\begin{theorem}\label{h3:thm:height-one-maps}
The integral classes \(\zeta\) and \(\delta\) above
and the canonical unit define an equivalence
\[
 (1,\zeta,\delta,\zeta\delta):
 R_1\vee\Sigma^{-1}R_1\vee
       \Sigma^{-3}R_1\vee\Sigma^{-4}R_1
      \xrightarrow{\ \simeq\ }L_{K(1)}T.
\]
The maps use the \(R_1\)-action and the multiplication
induced from \(T\).  In particular the first two
components are the \(K(1)\)-localizations of the
same maps from \(S\) used in
\Cref{h3:thm:height-two-test}.
\end{theorem}

\begin{proof}
The target is a \(K(1)\)-local algebra and hence a
module over the local unit \(R_1\).  Each homotopy
class displayed defines an actual map from the
indicated suspension of \(R_1\).  The product
\(\zeta\delta\) is formed using this algebra
multiplication, so no independent choice of a
degree-minus-four line is being made.

Modulo \(p\), the source groups are the direct sum
of the four shifts of
\[
 \pi_*(R_1/p)
     =\F_p[v_1^{\pm1}]\otimes\Lambda(h).
\]
By \Cref{h3:des:unit-directions}, the images of
\(1,\zeta,d,\zeta d\), together with their
\(h\)-multiples, are exactly the eight exterior
basis elements in the target coefficient
calculation.  The class \(\delta\) reduces to \(d\)
by \Cref{h3:prop:integral-delta}.
Products are detected by the cup action, so these
statements specify the leading filtered class of
each of the eight actual homotopy images.
The spectral sequence has collapsed with finite filtration.  In each
stem, filter the source homotopy group by assigning these eight
periodic basis elements the cohomological degrees of their displayed
leading classes.  The map is filtered, and its associated-graded map
is an isomorphism in every degree.  Induction on the finite filtration,
using the short exact sequences for successive filtration quotients,
shows that the original map of homotopy groups is an isomorphism.
This argument does not require an a priori splitting of the target
filtration.  The displayed map is therefore an isomorphism on
mod-\(p\) homotopy in every stem.

Its cofiber is \(K(1)\)-local, hence derived
\(p\)-complete by \Cref{h3:lem:height-one-completion}.  The vanishing of its mod-\(p\)
quotient makes multiplication by \(p\) an
equivalence, so all its quotients by \(p^\ell\)
vanish.  Derived completeness then makes the
cofiber zero.  This proves the equivalence,
including the asserted compatibility of its
first two components.
\end{proof}
